\section{实验设计}
\label{sec:experimental-setup}

本章给出实验环境、场景设置、对比算法、评价指标和统计检验方案。所有正式实验均使用第 \ref{sec:problem-formulation} 节定义的 B 样条控制点编码、同一目标函数、同一约束违反量和同一可行性优先接收准则。除明确说明的消融版本外，对比算法不得调用 CVF 场、状态自适应 CVF 权重、稀疏可行性保持或 CVF-AE 专属初始化机制。本文仅报告已有正式三维 UAV 路径规划实验，不引入额外未完成验证的通用约束基准或真实地图案例。

\subsection{实验环境}

实验在同一台本地工作站上完成。硬件和软件环境如下：CPU 为 Intel(R) Core(TM) i5-14600KF，主频 \SI{3.50}{GHz}；内存为 \SI{32}{GB}；操作系统为 64 位系统；MATLAB 版本为 R2022b。所有正式运行均按单进程串行方式执行，未启用并行计算。由于 wall-clock time 会受到 MATLAB 会话状态、repair 次数和批处理负载影响，本文不将绝对运行时间作为主要效率结论；计算开销主要由函数评价次数、CVF 触发规模和同批次评价比值刻画。

\subsection{测试场景}

三维规划空间在水平方向均为 $[0,100]\times[0,100]$，起点和终点分别为 $\mathbf{s}=[5,5,12]^\mathsf{T}$ 和 $\mathbf{g}=[95,95,18]^\mathsf{T}$。论文统一采用 Scene 1、Scene 2 和 Scene 3 三个场景编号。三个场景均包含建筑物障碍、柱状禁飞区和风险热点，但约束密度、高度上界和风险分布不同，用于检验算法在由普通城市走廊到更强低空约束走廊中的稳定性。

\begin{table}[tbp]
\centering
\caption{三维低空路径规划测试场景设置}
\label{tab:scene-setup}
\small
\begin{tabular}{@{}lcccc@{}}
\toprule
场景 & 障碍物数 & 禁飞区数 & 风险热点数 & 高度范围 \\
\midrule
Scene 1 & 6 & 3 & 3 & $[8,40]$ \\
Scene 2 & 10 & 2 & 3 & $[8,32]$ \\
Scene 3 & 17 & 3 & 5 & $[8,26]$ \\
\bottomrule
\end{tabular}
\end{table}

Scene 1 对应基准城市走廊，用于验证常规约束下的路径质量和可行性；Scene 2 对应中等强度约束走廊，用于考察更密集建筑和收紧高度上界下的可行路径形成能力；Scene 3 对应低空配送走廊，用于测试窄通道、风险热点和更严格高度上界共同作用下的鲁棒性。
Scene 1 使用参考高度 $z_{\mathrm{ref}}=16$；Scene 2 保持同一参考高度但将最大高度降至 $32$；Scene 3 将参考高度设为 $14$，最大高度进一步降至 $26$，并提高高度偏好项权重。该设置避免算法通过过高爬升绕开低空约束，使实验更贴近强约束低空飞行任务。

\subsection{公共参数与评价指标}

所有群智能算法采用相同的路径编码和评价预算。每条候选路径由 $K=6$ 个内部控制点表示，决策维度为 $18$，三次 B 样条路径采样点数为 $M=220$。正式主对比、近年改进算法补充实验和正式消融实验均采用 $N=30$ 的种群规模、$T=300$ 的最大迭代次数以及 $30$ 次独立运行。经典与主方法对比使用基础随机种子 20260630；近年改进算法补充实验使用基础随机种子 20260701，并与主对比结果按相同场景、相同预算和相同评价函数合并。参数稳健性检查只在 Scene 2 和 Scene 3 上进行，用于验证默认 CVF 参数是否具有稳定性，而不作为重新调参依据。

\begin{table}[tbp]
\centering
\caption{公共路径编码与评价参数设置}
\label{tab:common-setup}
\small
\begin{tabularx}{\linewidth}{@{}l X@{}}
\toprule
类别 & 设置 \\
\midrule
路径编码 & $K=6$，维度 $18$ \\
B 样条 & 三次 B 样条，$M=220$ \\
搜索预算 & $N=30$，$T=300$，30 次独立运行 \\
目标权重 & $w_L=1.00$，$w_E=0.60$，$w_R=1.20$，$w_S=0.15$，$w_H=0.80$，$w_B=0.80$ \\
违反惩罚 & $\lambda_{\mathrm{obs}}=150$，$\lambda_{\mathrm{nfz}}=200$，$\lambda_{\mathrm{curv}}=20$，$\lambda_{\mathrm{alt}}=30$ \\
运动约束 & $\theta_{\max}=55^\circ$，默认高度 $[8,40]$ \\
\bottomrule
\end{tabularx}
\end{table}

表 \ref{tab:common-setup} 仅保留复现实验所需的核心公共参数。路径编码只优化内部控制点，起点和终点固定；所有目标项和违反量均基于同一 B 样条采样路径评价；主对比、补充基线和消融实验保持相同搜索预算。惩罚项用于形成综合目标值，可行性比较仍以总违反量 $V$ 和 Deb 可行性优先规则为准；Scene 2 和 Scene 3 在此基础上覆盖更低高度上界及相应高度偏好设置。
主要评价指标包括：最终最佳适应度均值和标准差、可行率、最终违反量、首次获得可行解的迭代数、函数评价次数、平均排名以及轨迹合理性检查结果。对 CVF-AE 及其消融版本，额外统计 CVF 触发次数、CVF 候选成功次数、稀疏修复次数和修复成功次数。由于标量适应度不能完全反映边界贴靠和高风险绕行等轨迹问题，结果分析中将适应度、可行率、最终违反量和轨迹检查联合解释。

轨迹合理性检查采用固定规则生成，不根据算法结果临时调整。设 $z_{\max}$ 为场景高度上界，$z_{\mathrm{ref}}$ 为参考高度。高空阈值和极高空阈值定义为
\begin{equation}
    \begin{aligned}
    z_{\mathrm{high}}&=\min(z_{\max}-2,\ z_{\mathrm{ref}}+8),\\
    z_{\mathrm{vhigh}}&=\min(z_{\max}-1,\ z_{\mathrm{ref}}+16).
    \end{aligned}
    \label{eq:height-review-thresholds}
\end{equation}
对应比例为
\begin{equation}
    \begin{aligned}
    r_{\mathrm{high}}&=
    \frac{1}{M}\sum_{m=1}^{M}
    \mathbb I(z_m\geq z_{\mathrm{high}}),\\
    r_{\mathrm{vhigh}}&=
    \frac{1}{M}\sum_{m=1}^{M}
    \mathbb I(z_m\geq z_{\mathrm{vhigh}}).
    \end{aligned}
    \label{eq:height-review-rates}
\end{equation}
若 $r_{\mathrm{high}}>0.50$ 或 $r_{\mathrm{vhigh}}>0.20$，则标记为高空绕行风险。边界贴靠检查只统计去除首尾 $5\%$ 采样点后的路径内部段。令
\begin{equation}
    \begin{aligned}
    d_b(\mathbf p_m)=\min\{&
    x_m-x_{\min},\ x_{\max}-x_m,\\
    &y_m-y_{\min},\ y_{\max}-y_m\},
    \end{aligned}
    \label{eq:boundary-review-distance}
\end{equation}
$b_{\mathrm{review}}=\min(m_b,8)$，若内部采样点中 $d_b(\mathbf p_m)\leq b_{\mathrm{review}}$ 的比例超过 $0.65$ 且平均边界距离小于 $2b_{\mathrm{review}}$，则标记为边界贴靠风险。绕行比例定义为
\begin{equation}
    r_{\mathrm{detour}}=
    \frac{L}{\|\mathbf g-\mathbf s\|_2},
    \label{eq:detour-ratio}
\end{equation}
若 $r_{\mathrm{detour}}>1.80$，则标记为过度绕行。内部控制点若有超过 $65\%$ 位于边界检查带内，则标记为控制点贴边。综合复核标记由高空绕行、边界贴靠、过度绕行和控制点贴边四类标记取逻辑或得到。

\subsection{对比算法与消融版本}

主对比覆盖 11 个算法。本文方法为 CVF-AE；基础骨架基线为 AE；经典和常用群智能基线包括 PSO、GWO、WOA、HHO、DBO 和 CPO；近年改进算法补充组包括 GDSAO、ERIME 和 MSCSO。所有基线均使用同一 UAV 路径编码、同一目标函数、同一边界投影和同一可行性优先准则。新增的近年改进算法只接入其核心搜索机制，不调用本文方法的 CVF 场、约束状态调度或稀疏保持机制。

为避免基线被本文专属机制增强，正式主对比中基线算法不使用 CVF-AE 的约束感知初始化、CVF 候选分支或稀疏修复复用；它们仅共享公共的边界投影、B 样条解码、目标函数和 Deb 可行性优先比较。各算法参数采用本地统一实现中的固定配置，不针对单个场景进行事后调参。该设置使主对比反映完整 CVF-AE 框架相对于未嵌入本文专属约束引导机制的基线搜索器的性能差异；各模块的具体贡献由消融实验进一步解释。

消融实验用于识别 CVF-AE 各模块的作用，包含完整 CVF-AE、基础 AE、去除约束感知初始化版本、去除状态自适应 CVF 版本、去除稀疏保持版本和去除 CVF 场版本。消融实验与主对比使用相同的三场景、相同运行次数、相同种群规模和相同迭代次数。这样可以将性能差异主要归因于方法模块，而不是评价预算或场景设置变化。

\subsection{统计检验}

由于正式主对比、近年改进算法补充实验和消融实验中的算法随机种子均包含算法序号偏移，不同算法之间不能视为严格配对样本。因此，本文对 CVF-AE 与各基线的逐场景成对比较采用双侧 Wilcoxon rank-sum 检验，也即 Mann--Whitney U 检验 \citep{Mann1947RankSum}。每个场景内的多重成对比较采用 Holm 校正，显著性水平设为 $\alpha=0.05$ \citep{Holm1979Sequential}。同时，使用 Kruskal--Wallis 检验判断同一场景下多算法总体分布是否存在显著差异 \citep{Kruskal1952RankVariance}，并使用面向最小化问题定义的 Cliff's delta 报告效应量 \citep{Cliff1993Dominance}；正值表示 CVF-AE 的适应度分布更优。

统计结论仅用于说明最终最佳适应度分布差异，不单独决定工程优劣。对于可行率不足、最终违反量较高、边界贴靠明显或轨迹合理性检查风险较高的算法，即使其标量适应度较低，也需要在第 \ref{sec:results-discussion} 节中结合轨迹和约束指标进行讨论。
